\documentclass[11pt,twocolumn]{article}
\usepackage[margin=0.85in]{geometry}
\usepackage{amsmath,amssymb}
\usepackage{booktabs}
\usepackage{hyperref}
\usepackage{enumitem}
\usepackage{titlesec}

\titleformat{\section}{\large\bfseries}{\thesection}{1em}{}
\titleformat{\subsection}{\normalsize\bfseries}{\thesubsection}{1em}{}

\title{Gated Ensemble Neoantigen Prioritization:\\Variant Allele Expression as a Hard Filter\\Achieves 0.505 Recall@20 on 1.8M Peptides}
\author{Atris Research}
\date{March 2026}

\begin{document}
\maketitle

\begin{abstract}
We present a two-stage pipeline for neoantigen vaccine candidate selection
that achieves recall@20 of 0.505 in leave-one-patient-out cross-validation
across 99 patients and 1,787,710 peptides from the NeoRanking, TESLA, and
HiTIDE datasets. Our key insight is a hard gate on variant allele RNA-seq
support (alt\_support $> 0$) combined with MHC binding rank ($< 2.0$), which
eliminates 98.2\% of candidates while retaining 80.9\% of immunogenic
peptides (13.6$\times$ enrichment). A stacked GBT+RF ensemble applied to
gated candidates achieves an 83\% improvement over binding-rank-only
prioritization ($p = 0.0002$). We additionally identify TP53 R175H and
KRAS G12D as confirmed shared neoantigens immunogenic across independent
patients, and explain the BRAF V600E immunogenicity paradox. All experiments
use deterministic LOPO cross-validation with no data leakage.
\end{abstract}

\section{Introduction}

Personalized cancer vaccines targeting neoantigens---mutant peptides presented
on tumor cell surfaces---represent a promising immunotherapy approach.
The central computational challenge is selecting which of thousands of
candidate peptides to include in a vaccine, given that typically fewer than
0.01\% are immunogenic.

Current approaches rank candidates primarily by predicted MHC binding
affinity. We show that a simple two-gate filter based on variant allele
expression and binding rank, followed by an ensemble machine learning
scorer, substantially outperforms binding-only ranking.

\section{Data}

We use the NeoRanking benchmark (Bjerregaard et al., 2020), which
consolidates three datasets:

\begin{itemize}[nosep]
\item \textbf{NCI}: 1,377,707 peptides, 103 positive, 73 patients
\item \textbf{TESLA}: 301,741 peptides, 34 positive, 16 patients
\item \textbf{HiTIDE}: 108,262 peptides, 41 positive, 10 patients
\end{itemize}

Total: 1,787,710 peptides, 178 confirmed CD8+ immunogenic across 99 patients
and 14 cancer types. Class imbalance: 99.99\% negative.

\section{Methods}

\subsection{Two-Gate Filter}

We apply two hard gates sequentially:

\begin{enumerate}[nosep]
\item \textbf{Variant allele support $> 0$}: The mutation must have RNA-seq
reads supporting the alternate allele. This confirms the mutation is
transcribed from the tumor genome.
\item \textbf{MHC binding rank $< 2.0$}: The mutant peptide must bind
to at least one patient HLA allele with rank percentile below 2\%.
\end{enumerate}

\subsection{Ensemble Scoring}

Gated candidates are scored by a stacked ensemble:
\begin{itemize}[nosep]
\item Gradient Boosted Trees (60\% weight): 100 trees, depth 3
\item Random Forest (40\% weight): 100 trees, depth 5
\end{itemize}

Features (26 total): binding ranks (MixMHC, NetMHCpan, PRIME), stability
rank, TAP score, proteasomal cleavage, expression (RNA-seq TPM, alt support,
GTEx, TCGA), differential agretopicity, wildtype similarity, CSCAPE driver
score, cancer cell fraction, peptide length, and interaction terms
(binding$\times$expression, log transforms).

\subsection{Evaluation}

Leave-one-patient-out (LOPO) cross-validation: for each of 73 patients
with $\geq 1$ positive peptide, train on all other patients and evaluate
recall@$k$ (fraction of immunogenic peptides in the top $k$ predictions).

\section{Results}

\subsection{Gate Performance}

\begin{table}[h]
\centering
\small
\begin{tabular}{lrrr}
\toprule
\textbf{Gate} & \textbf{Peptides} & \textbf{Pos Kept} & \textbf{Rate} \\
\midrule
None & 1,787,710 & 100\% & 0.010\% \\
Alt $> 0$ & 536,857 & 85.4\% & 0.028\% \\
Alt $> 0$ + Bind $< 2$ & 32,669 & 80.9\% & 0.44\% \\
+ Expr $> 10$ TPM & 22,040 & 74.7\% & 0.60\% \\
\bottomrule
\end{tabular}
\caption{Gate filter performance. The two-gate filter eliminates 98.2\% of
candidates with 80.9\% positive retention.}
\end{table}

\subsection{Recall@20 Comparison}

\begin{table}[h]
\centering
\small
\begin{tabular}{lrrl}
\toprule
\textbf{Method} & \textbf{R@20} & \textbf{SE} & \textbf{$p$ vs bind} \\
\midrule
Gated GBT+RF & \textbf{0.505} & 0.051 & 0.0002 \\
Gated GBT & 0.497 & 0.051 & 0.0002 \\
Ungated GBT & 0.432 & 0.051 & 0.003 \\
Binding only & 0.276 & 0.046 & --- \\
\bottomrule
\end{tabular}
\caption{LOPO recall@20. Gated stacked ensemble achieves 83\% improvement
over binding-only baseline.}
\end{table}

\subsection{Recall@$k$ Curve}

At $k=50$ candidates, gated GBT achieves 0.649 recall; at $k=100$, 0.751.
For a vaccine synthesizing 20 peptides, we recover half of all immunogenic
candidates per patient.

\subsection{Variant Allele Support Is the Key Feature}

The \texttt{rnaseq\_alt\_support} feature shows the strongest separation
between positive and negative peptides: median 39 reads (positive) vs.\
0 (negative), Mann-Whitney $p = 2.6 \times 10^{-57}$. This confirms that
immunogenic neoantigens must be actively transcribed from the mutant allele.

Feature importance (GBT): binding rank (netMHCpan) 19.7\%, alt support
12.5\%, TCGA expression 8.0\%, TAP score 5.6\%, binding score 5.2\%.

\subsection{Expression $\times$ Binding Interaction}

Strong binding + high expression yields 21.5$\times$ enrichment over
weak binding + low expression. Neither feature alone is sufficient:
strong binding alone gives 1.8$\times$, high expression alone gives
0.2$\times$. The ensemble captures this non-linear interaction.

\subsection{Cross-Dataset Transfer}

Training on NCI+HiTIDE and evaluating on TESLA: AUC = 0.995.
Training on TESLA and evaluating on NCI: AUC = 0.959.
Positive peptide profiles are statistically indistinguishable between
datasets for binding ($p = 0.52$) and expression ($p = 0.93$).

\subsection{Error Analysis}

Of 178 positives: 73 found in top 20 (hits), 71 passed gate but ranked
$>$20 (misses), 34 failed the gate entirely. Gate failures: 26 lack
alt support (no RNA-seq evidence), 8 have weak binding ($\geq 2.0$).
Hits have median binding 0.05 vs.\ 0.20 for misses and stability
0.30 vs.\ 1.10.

\section{Shared Neoantigens}

\subsection{Confirmed Shared Neoantigens}

Two mutations are immunogenic in $\geq 2$ independent patients:
\begin{itemize}[nosep]
\item \textbf{TP53 R175H}: 2/6 patients immunogenic (33\%), presented by
HLA-A*02:01 (29\% of population)
\item \textbf{KRAS G12D}: 2/10 patients immunogenic (20\%), presented by
12+ HLA alleles ($>$50\% population coverage)
\end{itemize}

\subsection{The BRAF V600E Paradox}

BRAF V600E appears in 20 patients (most common driver) but is 0\%
immunogenic. We show this is because only 5.4\% of BRAF V600E peptides
bind strongly to any HLA allele (median rank 14.0). BRAF patients are
actually \textit{more} immunogenic than non-BRAF patients (0.0175\% vs.\
0.0084\%) for \textit{other} mutations. BRAF V600E should be targeted
pharmacologically, not immunologically.

\section{Additional Findings}

\begin{itemize}[nosep]
\item \textbf{Mutation burden}: significantly correlates with immunogenicity
(Spearman $r = 0.226$, $p = 0.025$)
\item \textbf{HLA alleles}: A*68:01 (7.4$\times$) and B*35:03
(6.3$\times$) are most enriched for immunogenicity
\item \textbf{AA substitutions}: charge-changing mutations (P$\to$R,
R$\to$M) are 5$\times$ more immunogenic than average
\item \textbf{Driver genes}: 3.7$\times$ enriched (17.4\% of positives
vs.\ 4.7\% of negatives)
\item \textbf{Clonality}: higher CCF correlates with immunogenicity
($p = 0.0007$)
\item \textbf{9-mers}: 15$\times$ more immunogenic than 8-mers
\item \textbf{Simple scores fail}: no formula without ML replicates
the ensemble's improvement ($p = 0.003$)
\item \textbf{Patient prediction fails}: LOOCV AUC 0.420; 73.7\%
of patients have $\geq 1$ positive, so assume universal vaccination
\end{itemize}

\section{Clinical Protocol}

Based on these findings, we propose:
\begin{enumerate}[nosep]
\item Sequence tumor (WGS/WES) + RNA-seq
\item Call neoantigens with patient-matched HLA alleles
\item \textbf{Gate}: alt support $> 0$ AND binding rank $< 2.0$
\item \textbf{Score}: GBT+RF ensemble on gated candidates
\item \textbf{Select}: top 20 for vaccine synthesis
\item Expected yield: $\sim$50\% of immunogenic peptides in 20 candidates
\end{enumerate}

Cost: WGS (\$200) + RNA-seq (\$50) + HLA typing (\$100) = \$350/patient.
Compute: $<$1 minute on a laptop. No RNA-seq? Use TCGA cancer-type
averages (same AUC, Paper 01).

\section{Conclusion}

A two-gate filter (variant allele expression + MHC binding) combined
with ensemble machine learning achieves 0.505 recall@20 for neoantigen
prioritization---an 83\% improvement over the binding-only standard of care.
The gate alone provides 13.6$\times$ enrichment while eliminating 98.2\%
of candidates. These results are validated across three independent
datasets with consistent cross-dataset transfer.

\end{document}
